\reading{LTR-2}{Liquidity and Leverage}{DEFER}{0}

\begin{lrkeybox}[Learning objectives — official 2026 spine wording]
\begin{enumerate}[label=\textbf{\alph*.}]
\item Differentiate between sources of liquidity risk and describe specific challenges faced by different types of financial institutions in managing liquidity risk.
\item Summarize the asset-liability management process at a fractional reserve bank, including the process of liquidity transformation.
\item Compare transactions used in the collateral market and explain risks that can arise through collateral market transactions.
\item Describe the relationship between leverage and a firm's return profile (including the leverage effect), and explain the impact of different types of transactions on a firm's leverage and balance sheet.
\item Describe and compare methods to measure and manage funding liquidity risk and transactions liquidity risk.
\item Calculate the expected transactions cost and the spread risk factor for a transaction and calculate the liquidity adjustment to VaR for a position to be liquidated over a number of trading days.
\item Discuss interactions between different types of liquidity risk and explain how liquidity risk events can increase systemic risk.
\end{enumerate}
{\footnotesize\color{FRMGrey}Source: Allan Malz, \textit{Financial Risk Management: Models, History, and Institutions}, Chapter 12.}
\end{lrkeybox}

\begin{lrtrapbox}[A numbering warning specific to this reading]
Both source layers label this reading's objectives, and they disagree. The Crash layer calls them \texttt{LO 65.a--g}; the Class Lecture layer calls them \texttt{LO 60.a--g}. Worse, the Crash layer's \texttt{LO 65} prefix is reused by the Lecture layer for \textbf{LTR-7 Monitoring Liquidity}, so \texttt{65.b} means two different objectives depending on which layer you are holding. Everything here is renumbered to the spine. Retire the old labels.
\end{lrtrapbox}

% =====================================================================
\lo{LTR-2 a}{Differentiate between sources of liquidity risk and describe specific challenges faced by different types of financial institutions in managing liquidity risk.}

\begin{lrdefbox}[Liquid asset, liquid market]
An asset is liquid if it is \emph{close to cash} — it can be sold quickly, cheaply, and without moving the price too much. A market is liquid if positions can be unwound quickly, cheaply (that is, at low transactions costs), and without undue price deterioration.
\end{lrdefbox}

\subsection{The two main types of liquidity risk}

\begin{center}
\small
\begin{tabularx}{\textwidth}{@{}p{3.8cm} X@{}}
\toprule
\rowcolor{FRMSoft}\textbf{Type} & \textbf{What it is} \\
\midrule
\textbf{1. Transactions liquidity risk} & The risk that buying or selling an asset will cause its price to move unfavourably. It pertains to \textbf{assets and financial markets}. \\
\addlinespace
\textbf{2. Funding liquidity risk} & Arises when a borrower's creditworthiness deteriorates, or when market conditions worsen. Creditors may withdraw credit or change terms, potentially forcing the borrower to unwind positions or face losses. \\
\bottomrule
\end{tabularx}
\end{center}

\begin{lrtrapbox}[Role: which side of the balance sheet]
Transactions liquidity is a property of the \textbf{asset and its market}. Funding liquidity is a property of the firm's \textbf{liabilities and its creditors' willingness to lend}. GARP tests this as a straight definitional discriminator and it is named in \S5.4 as one of the standing confusions.
\end{lrtrapbox}

\subsection{Three structural features that generate the risk}

\begin{itemize}
\item \term{Maturity mismatch}: borrowing short term to finance investments that require a longer time to become profitable. Intermediaries engage in maturity mismatch because it is generally profitable.
\item Funding long-term assets with short-term debt exposes an intermediary to \term{rollover risk} or even \term{cliff risk}.
\item \term{Systemic risk} is the risk that the overall financial system is impaired due to severe financial stress.
\end{itemize}

% =====================================================================
\lo{LTR-2 b}{Summarize the asset-liability management process at a fractional reserve bank, including the process of liquidity transformation.}

\subsection{Liquidity transformation by banks}

Commercial banks have liabilities that are typically shorter-term and more liquid than their assets. Deposits constitute a significant portion, approximately 60\%, of these liabilities.

\begin{itemize}
\item \term{Fractional-reserve banking.} Most banks historically operate on a fractional-reserve system: they keep only a fraction of deposited funds in liquid assets and use the rest to make loans. This approach is known as \term{asset-liability management (ALM)}.
  \begin{itemize}
  \item If a bank receives \$100 in deposits, it might keep \$10 for potential withdrawals and lend out the remaining \$90.
  \item Alternatively, banks can use owners' equity or funds raised from capital markets to make loans, while keeping cash or highly liquid assets equal to their deposits in reserve.
  \end{itemize}
\item \term{Asset-liability management (ALM)} is the practice of balancing a bank's assets (loans) with its liabilities (deposits) to manage liquidity risk.
\item \term{Sources of funding for banks.} Deposits are the primary source, but banks also use wholesale funding such as commercial paper and bonds.
\end{itemize}

\subsection{Fragility of commercial banking}

\begin{itemize}
\item In a fractional-reserve system, if depositors wish to make withdrawals in excess of a bank's reserves, and the bank cannot liquidate enough loans or other assets to meet the demand immediately, it is forced into \term{suspension of convertibility}.
\item At the extreme, all or a large number of depositors may ask for the return of their money simultaneously — a \term{bank run}.
\item Apart from deposits, banks are generally dependent on short-term financing, exposing them to \term{rollover risk}: the risk that short-term debt cannot be refinanced, or can be refinanced only on highly disadvantageous terms. (Example: LTCM.)
\end{itemize}

\subsection{Structured credit products and the off-balance-sheet vehicles}

\begin{itemize}
\item \term{Structured credit products (SCPs)}, such as ABSs and MBSs, pool assets like mortgages and make them tradable, matching investor needs with assets and avoiding funding liquidity issues through maturity matching. They offer maturity transformation, but create liquidity risk if investor funding is short-term.
\item \textbf{Liquidity risk from short-term financing.} Relying on short-term financing for structured credit products led to liquidity risk, a key factor in the 2007--2009 subprime crisis.
\item \textbf{Types of short-term financing.} Securities lending uses structured credit products as collateral for short-term loans. Off-balance-sheet vehicles, such as \term{special-purpose vehicles (SPVs)}, issue secured debt (for example ABCP) to finance assets, receiving liquidity and credit support.
\item \term{Structured investment vehicles (SIVs)} are similar to ABCP conduits but lack full liquidity and credit support.
\end{itemize}

\begin{lrkeybox}[Maturity and liquidity transformation off the balance sheet]
\begin{itemize}
\item Off-balance-sheet vehicles perform \textbf{maturity transformation} by holding longer-maturity assets than the short-term ABCP used to fund them.
\item They also perform \textbf{liquidity transformation} by creating more liquid, shorter-term ABCP than their underlying assets.
\item Despite being off balance sheet, these vehicles did \textbf{not} entirely transfer risk.
\item They contributed to financial system fragility during the subprime crisis by amplifying leverage when asset values declined.
\end{itemize}
\end{lrkeybox}

\subsection{Systematic funding liquidity risk}

During the subprime crisis, systematic funding risks affected various market sectors. Shorter-term loans increased liquidity risks as borrowers needed to refinance to repay short-term obligations. Liquidity issues arose for a variety of investment strategies:

\begin{enumerate}
\item \textbf{Leveraged buyouts (LBOs).} LBOs and private equity relied heavily on leveraged loans, a dominant form of syndicated bank loan. CLOs and CDOs funded LBOs. As funding dried up, LBO deals fell apart, leading to significant losses for banks.
\item \textbf{Merger arbitrage hedge funds.} Financing difficulties caused some announced mergers to be abandoned, producing losses for funds running the strategy.
\item \textbf{Convertible arbitrage hedge funds.} The strategy relies on leverage extended by broker-dealers. Funding unavailability caused sharp drops in convertible bond values; market liquidity problems and investor redemptions made it worse. Limited investor interest made selling difficult, contributing to price declines. Despite widening price gaps, arbitrage capital did not re-enter the market.
\end{enumerate}

\subsection{Money market mutual funds}

\begin{enumerate}[label=\alph*)]
\item MMMFs are investment vehicles similar to banks, allowing investors to write cheques and make electronic transfers while promising to repay on demand. They invest primarily in high credit quality, short-term instruments with maturities from a few weeks to a few months.
\item Despite holding relatively safe assets, the values of underlying assets can change, potentially limiting redemptions if asset values decline.
\item MMMF liabilities are more liquid than their investments, resembling the liquidity structure of banks. MMMFs use the \term{amortized cost method} under the SEC's Rule 2a-7, so their assets are not marked to market daily — very short-term securities are less likely to revalue on interest rate and credit spread changes.
\item MMMFs typically maintain a notional share value of \$1.00, but credit write-downs or liquidity issues can push net asset value below \$1.00, a situation known as \term{breaking the buck}.
\item Liquidity risk, including investor runs in adverse conditions, can force MMMFs to sell assets at a loss, leading to write-downs and potential instances of breaking the buck.
\end{enumerate}

% =====================================================================
\lo{LTR-2 c}{Compare transactions used in the collateral market and explain risks that can arise through collateral market transactions.}

Collateral markets serve two key roles: facilitating borrowing using assets (including securities), and enabling short positions in securities. They work by using securities as collateral for loans, making lenders more willing to lend at lower interest rates.

\begin{lrkeybox}[Four features of the collateral market]
\begin{itemize}
\item \textbf{Enhanced borrowing.} Firms prefer lending at lower rates when loans are secured with collateral. Collateralized loans vary in duration, including short-term and overnight loans with automatic extensions.
\item \textbf{Haircuts and variation margin.} Collateralized loans involve lending less than the full collateral value; the difference is the \term{haircut}. To maintain full collateralization as collateral values fluctuate, \term{variation margin} is paid, protecting lenders.
\item \textbf{Rehypothecation.} Collateral assets are often reloaned or pledged again, allowing collateral to circulate.
\item \textbf{Modern finance role.} Collateral's importance has grown with securitization, which creates securities that can be used as collateral for other transactions. Life insurers use high-quality asset portfolios as collateral for cost-effective borrowing and reinvestment; hedge funds pledge securities to finance portfolios at better rates than unsecured loans.
\end{itemize}
\end{lrkeybox}

\subsection{The four forms a collateral market takes}

\begin{center}
\small
\begin{tabularx}{\textwidth}{@{}p{3.1cm} X X@{}}
\toprule
\rowcolor{FRMSoft}\textbf{Transaction} & \textbf{Mechanics and purpose} & \textbf{Major risks} \\
\midrule
\textbf{1. Margin loans} &
Lending to finance a security transaction, collateralized by that security. The simplest form of collateral market, used mainly by investors taking leveraged \textbf{long} positions, most often in equities. &
\textbf{Repledge}: the broker is likely to use the customer's collateral to borrow in the secured money market to obtain the funds it lends to margin customers; collateral may also be repledged to borrow another security rather than cash. \textbf{Margin call}: if the market value of the long position declines, the broker loses collateral protection and calls the customer. \\
\addlinespace
\textbf{2. Total return swaps (TRS)} &
One party pays a fixed fee in exchange for the total return (income \emph{and} capital gains) on a reference asset, typically a stock. The fee payer earns the asset's return without owning it. &
The party providing the return (for instance a hedge fund) is, in essence, \textbf{short} the asset. \\
\addlinespace
\textbf{3. Repurchase agreements} &
Matched pairs of a spot sale and a forward repurchase of a security; collateralization is achieved by selling the security temporarily to the lender. \textbf{Reverse repo} is often used to finance long positions in securities, typically bonds; \textbf{repo} is usually used by bond owners to borrow cash. &
Repo has expanded from Treasuries to high-yield bonds, whole loans and more recently structured credit products. Mechanics resemble margin loans, and like margin lending, repo creates a straightforward liability on the economic balance sheet. \\
\addlinespace
\textbf{4. Securities lending} &
One party lends a security to another for a fee, generally called a \term{rebate}. The security lender continues to receive dividend and interest cash flows from the security. &
Two typical patterns: equities held at a custodian or prime broker in \emph{street name} so they can be rehypothecated to a short seller; and fixed-income lending aimed at earning a spread between less- and more-risky bonds. \\
\bottomrule
\end{tabularx}
\end{center}
\figcap{Read the middle column for direction. Margin loans finance a \textbf{long}; securities lending exists to establish a \textbf{short}; in a TRS the return \emph{provider} is effectively short. A four-statement audit on this reading almost always turns on which party ends up long and which short.}

\begin{lrkeybox}[Where the focus has shifted]
Securities lending has typically focused on the \emph{securities} rather than the cash collateral, usually to establish short positions. In recent years the focus of their use has shifted toward borrowing \emph{cash} collateral.
\end{lrkeybox}

% =====================================================================
\lo{LTR-2 d}{Describe the relationship between leverage and a firm's return profile (including the leverage effect), and explain the impact of different types of transactions on a firm's leverage and balance sheet.}

\begin{lrdefbox}[The leverage effect]
Leverage matters because it provides an opportunity to increase returns to equity investors. The \term{leverage effect} is the increase in equity returns that results from increasing leverage, and is equal to the difference between the return on assets and the cost of funding.
\end{lrdefbox}

\begin{lrfmlbox}[Return on equity under leverage]
\[ R_e \;=\; L \times R_a \;-\; (L-1) \times R_d \]
\vk{$R_e$ = return on equity; $R_a$ = return on assets; $R_d$ = cost of debt funding; $L$ = leverage ratio (assets divided by equity). At $L=1$ the firm is unlevered and $R_e = R_a$.}
\end{lrfmlbox}

Rearranged, the identity says something the formula hides: $R_e = R_a + (L-1)(R_a - R_d)$. Equity return is the asset return plus $(L-1)$ times the \emph{spread} between what the assets earn and what the funding costs. The leverage effect is that spread term.

\begin{center}
\begin{tikzpicture}
\begin{axis}[
  width=12.5cm, height=6.6cm,
  xlabel={Leverage ratio $L$ (assets / equity)},
  ylabel={Return on equity $R_e$ (\%)},
  xlabel style={font=\small\sffamily}, ylabel style={font=\small\sffamily},
  tick label style={font=\scriptsize\sffamily},
  xmin=1, xmax=10, ymin=-16, ymax=34,
  ymajorgrids, grid style={FRMRule, dashed},
  legend style={font=\scriptsize\sffamily, draw=FRMRule, fill=white,
    at={(0.5,-0.26)}, anchor=north, legend columns=-1, column sep=10pt},
]
\addplot[FRMGreen, very thick, domain=1:10, samples=2]{8 + (x-1)*3};
\addlegendentry{$R_a = 8\%$, $R_d = 5\%$ \ (assets out-earn funding)}
\addplot[FRMRed, very thick, domain=1:10, samples=2]{3 - (x-1)*2};
\addlegendentry{$R_a = 3\%$, $R_d = 5\%$ \ (funding out-costs assets)}
\addplot[FRMNavy, dashed, domain=1:10, samples=2]{0};
\node[font=\scriptsize\sffamily, fill=white, inner sep=1.5pt, anchor=west] at (axis cs:6.1,27) {slope $= R_a - R_d = +3$};
\node[font=\scriptsize\sffamily, fill=white, inner sep=1.5pt, anchor=west] at (axis cs:5.9,-12) {slope $= R_a - R_d = -2$};
\node[font=\scriptsize\sffamily, fill=white, inner sep=1.5pt, anchor=west] at (axis cs:1.25,9.6) {unlevered: $R_e = R_a$};
\draw[-{Stealth[length=4pt]}, FRMNavy] (axis cs:1.22,9.2) -- (axis cs:1.05,8.2);
\end{axis}
\end{tikzpicture}
\figcap{The identity $R_e = L R_a - (L-1) R_d$ plotted for two funding spreads. Leverage is a multiplier on the \emph{sign} of $R_a - R_d$, not a generator of return: the two lines pivot around the same unlevered point and diverge linearly. The red line crosses zero at $L = 2.5$ and keeps going — which is the whole of the leverage lesson in one intersection.}
\end{center}

\subsection{Transactions and leverage}

\begin{enumerate}[label=\alph*)]
\item \textbf{Margin loans}: the leverage ratio increases with borrowed funds used for stock purchases.
\item \textbf{Short positions}: higher leverage, because the entire value of the stock is borrowed.
\item \textbf{Derivatives}: implicit leverage from synthetic long or short positions, calculated using the cash-equivalent market value of the underlying assets.
\end{enumerate}

\begin{lrtrapbox}[Implicit leverage does not show up as debt]
A derivative position carries leverage without a borrowing line on the balance sheet. The measurement convention is the \textbf{cash-equivalent market value of the underlying}, not the premium paid or the margin posted. A distractor built from the margin posted is the obvious one to offer.
\end{lrtrapbox}

% =====================================================================
\lo{LTR-2 e}{Describe and compare methods to measure and manage funding liquidity risk and transactions liquidity risk.}

\begin{lrdefbox}[Transactions liquidity risk, by its sources]
Related to the costs of finding a counterparty, the institutions assisting in that search, and the costs of inducing a counterparty to hold a position.
\end{lrdefbox}

\subsection{The four components of transactions liquidity risk}

\begin{enumerate}[label=\alph*)]
\item \term{Trade processing costs}: finding a counterparty, processing, clearing, and settling trades.
\item \term{Inventory management}: dealers holding long or short inventories to provide trade immediacy, requiring compensation for volatility exposure.
\item \term{Adverse selection}: differentiating between liquidity traders and information traders, with compensation for that risk taken through the bid-ask spread.
\item \term{Differences of opinion}: difficulty in finding a counterparty when market participants agree on the state of the market.
\end{enumerate}

\begin{lrtrapbox}[Differences of opinion runs the opposite way to intuition]
The difficulty arises when participants \textbf{agree}, not when they disagree. Agreement means everyone wants the same side and nobody will take the other. This is a polarity trap sitting inside an item that looks purely definitional.
\end{lrtrapbox}

\subsection{Three characteristics used to measure market liquidity}

\begin{center}
\small
\begin{tabularx}{\textwidth}{@{}p{3cm} X@{}}
\toprule
\rowcolor{FRMSoft}\textbf{Characteristic} & \textbf{What it measures} \\
\midrule
\textbf{Tightness} (or width) & The cost of a round-trip transaction, measured by the bid-ask spread and brokers' commissions. \\
\textbf{Depth} & How large an order must be to move the price adversely. \\
\textbf{Resiliency} & The length of time it takes lumpy orders to move the market away from the equilibrium price. \\
\bottomrule
\end{tabularx}
\end{center}

\begin{lrkeybox}[Managing funding liquidity: the hedge fund case]
\begin{enumerate}[label=\alph*)]
\item Hedge funds manage liquidity via cash, unpledged assets, and unused borrowing capacity.
\item In times of market stress, redemption requests may require managers to unwind positions rapidly, exposing the fund to transactions liquidity risk. If this happens to many funds at once, fire sales may result.
\end{enumerate}
\end{lrkeybox}

% =====================================================================
\lo{LTR-2 f}{Calculate the expected transactions cost and the spread risk factor for a transaction and calculate the liquidity adjustment to VaR for a position to be liquidated over a number of trading days.}

\begin{lrfmlbox}[Expected transactions cost and the spread risk factor]
\[ \text{Expected transactions cost} \;=\; P \times \frac{s}{2} \]
\[ \text{99\% confidence interval} \;=\; \pm\, P \times \tfrac{1}{2}\left(s + 2.33\,\sigma_s\right) \]
\vk{$P$ = the value of the position; $s$ = the expected (mean) proportional bid-ask spread; $\sigma_s$ = the standard deviation of the proportional spread. The $2.33\sigma_s$ term is the \term{spread risk factor} — the allowance for the spread being wider than expected at 99\% confidence.}
\end{lrfmlbox}

\begin{lrfmlbox}[Liquidity adjustment to VaR over a multi-day liquidation]
\[ \text{LVaR} \;=\; \text{VaR} \times \sqrt{\frac{(1+T)\,(1+2T)}{6T}} \]
\vk{$T$ = the number of trading days required for the orderly liquidation of the position. The multiplier assumes the position is unwound in equal instalments over $T$ days, so the average exposure remaining is less than the full position but the exposure lasts longer.}
\end{lrfmlbox}

\begin{center}
\begin{tikzpicture}
\begin{axis}[
  width=12.5cm, height=6.4cm,
  xlabel={Liquidation horizon $T$ (trading days)},
  ylabel={LVaR $\div$ VaR},
  xlabel style={font=\small\sffamily}, ylabel style={font=\small\sffamily},
  tick label style={font=\scriptsize\sffamily},
  xmin=1, xmax=20, ymin=0.9, ymax=3.1,
  xtick={1,2,5,10,15,20},
  ymajorgrids, grid style={FRMRule, dashed},
]
\addplot[FRMNavy, very thick, domain=1:20, samples=80]{sqrt((1+x)*(1+2*x)/(6*x))};
\addplot[only marks, mark=*, mark size=1.8pt, FRMGold] coordinates {(1,1.000)(2,1.118)(5,1.483)(10,1.962)(20,2.679)};
\node[font=\scriptsize\sffamily, fill=white, inner sep=1.5pt, anchor=south west] at (axis cs:1.2,1.02) {$T{=}1 \Rightarrow 1.000$};
\node[font=\scriptsize\sffamily, fill=white, inner sep=1.5pt, anchor=south west] at (axis cs:4.6,1.53) {$T{=}5 \Rightarrow 1.483$};
\node[font=\scriptsize\sffamily, fill=white, inner sep=1.5pt, anchor=south east] at (axis cs:19.6,2.52) {$T{=}20 \Rightarrow 2.679$};
\draw[-{Stealth[length=4pt]}, FRMGold] (axis cs:18.6,2.52) -- (axis cs:19.7,2.66);
\end{axis}
\end{tikzpicture}
\figcap{The multiplier $\sqrt{(1+T)(1+2T)/(6T)}$ against the liquidation horizon. It equals exactly 1.000 at $T=1$, which is the sanity check to run before locking any answer: a one-day unwind is ordinary VaR. It grows roughly as $\sqrt{T/3}$ for large $T$ — slower than the naive $\sqrt{T}$, because the position shrinks as it is sold.}
\end{center}

\begin{lrtrapbox}[Three ways this calculation is corrupted]
\begin{itemize}
\item \textbf{Using $\sqrt{T}$ instead of the full multiplier.} At $T=10$, $\sqrt{T} = 3.162$ against the correct 1.962. That wrong value will be one of the four options.
\item \textbf{Confidence twin.} The spread risk factor uses \textbf{2.33} (99\%). Substituting 1.645 (95\%) produces a clean, plausible, wrong number.
\item \textbf{Halving.} Both the expected cost and the confidence interval carry the $\tfrac{1}{2}$. Dropping it doubles the answer — again, a value that will be offered.
\end{itemize}
\end{lrtrapbox}

% =====================================================================
\lo{LTR-2 g}{Discuss interactions between different types of liquidity risk and explain how liquidity risk events can increase systemic risk.}

\begin{itemize}
\item Liquidity risks in finance are \textbf{interconnected and can worsen each other}. If collateral requirements rise, a party may have to sell assets early, increasing both funding and transactions liquidity risk. Leverage is crucial here and can force asset sales, reducing the pool of potential buyers and causing asset values to drop, potentially leading to a \term{debt-deflation crisis}.
\item \textbf{Transactions liquidity can also affect funding liquidity.} When a hedge fund faces redemptions, selling highly liquid assets minimises price impact but leaves the fund with a less liquid portfolio. Conversely, selling illiquid assets can lead to realized losses, adding to funding liquidity pressure.
\item \textbf{Economy-wide liquidity levels impact systemic risk.} During market stress, liquidity tends to tighten exactly when it is needed most, potentially creating systemic risk events by disrupting payment and settlement systems and affecting many market participants simultaneously.
\end{itemize}

\needspace{14\baselineskip}
\begin{center}
\begin{tikzpicture}[
  node distance=6mm,
  bx/.style={draw=FRMNavy, fill=PaleNavy, rounded corners=1pt, align=center,
             font=\scriptsize\sffamily, text width=2.9cm, inner sep=4pt},
  bxr/.style={draw=FRMRed, fill=PaleRed, rounded corners=1pt, align=center,
             font=\scriptsize\sffamily\bfseries, text width=3.0cm, inner sep=4pt},
  ar/.style={-{Stealth[length=4.5pt]}, FRMGold, thick}
]
\node[bx] (a) at (0,0) {Collateral requirements rise / creditworthiness falls};
\node[bx] (b) at (4.6,0) {Funding liquidity risk: positions must be unwound early};
\node[bx] (c) at (9.2,0) {Forced sales into a thin market};
\node[bx] (d) at (9.2,-2.3) {Transactions liquidity risk: prices move against the seller};
\node[bx] (e) at (4.6,-2.3) {Asset values fall; remaining collateral is worth less};
\node[bxr] (f) at (0,-2.3) {Debt-deflation spiral $\rightarrow$ systemic risk};

\draw[ar] (a) -- (b);
\draw[ar] (b) -- (c);
\draw[ar] (c) -- (d);
\draw[ar] (d) -- (e);
\draw[ar] (e) -- (f);
\draw[ar] (f.north) -- ($(a.south)$);
\end{tikzpicture}
\figcap{The loop closes — that is the whole content of this objective. Funding pressure creates transactions pressure, transactions pressure destroys collateral value, and destroyed collateral value is fresh funding pressure. Leverage sets how fast the loop turns.}
\end{center}

\begin{lrtrapbox}[Consolidated trap summary — LTR-2]
\begin{itemize}
\item \textbf{Sibling.} Transactions liquidity = asset and market. Funding liquidity = liabilities and creditors. Named in \S5.4 as a standing confusion.
\item \textbf{Polarity.} \emph{Differences of opinion} is a source of illiquidity when participants \textbf{agree}, not when they disagree.
\item \textbf{Role.} Margin loan finances a long; securities lending enables a short; in a TRS the return \emph{provider} is short.
\item \textbf{Calculation.} LVaR multiplier is $\sqrt{(1+T)(1+2T)/(6T)}$, not $\sqrt{T}$. It equals 1.000 at $T=1$ — use that as the execution check.
\item \textbf{Confidence twin.} Spread risk factor uses 2.33, not 1.645.
\item \textbf{Sign.} The leverage effect amplifies $R_a - R_d$ including when it is \textbf{negative}. Leverage is a multiplier, not a return generator.
\item \textbf{Scope.} SIVs resemble ABCP conduits but \textbf{lack} full liquidity and credit support. Off-balance-sheet placement did not transfer the risk.
\item \textbf{Numbering.} Old labels \texttt{65.x} and \texttt{60.x} both refer to this reading in the source notes, and \texttt{65.x} also refers to LTR-7. Use spine IDs only.
\end{itemize}
\end{lrtrapbox}
